\subsection{Architectural Pitfalls of Argument Evaluation}

Default arguments are evaluated when the \texttt{def} statement executes, not on each call. The resulting object is stored on the function object. Immutable defaults are safe; mutable defaults persist and accumulate state across calls.

\subsubsection{The Mutable Default Arguments Trap: Why \texttt{def func(x=[])} Reuses One Persistent Mutable Object}

%<<<PROGRAM:programs/arc03>>>


One list object is created at definition time and reused whenever \texttt{lines} is omitted.

\subsubsection{Static Expression Evaluation: Why Default Values Are Created at Function Definition Time}

%<<<PROGRAM:programs/arc02>>>


\texttt{make\_default\_tag()} runs once during \texttt{def}, not on every call.

\subsubsection{Defending Against State Contamination Using the Immutable \texttt{None} Idiom and Sentinel Guard Clauses}

%<<<PROGRAM:programs/arc01>>>


Use \texttt{None} as the default sentinel; allocate mutable state inside the function body on each call.
